\documentclass{situationist-preprint}

\title{Meinong's Jungle Cleared}
\msid{MET-24-b}
\msfamily{MET-24}
\msdiscipline{Metaphysics}
\mstype{Companion manuscript}
\mskeywords{Meinong; nonexistent objects; incomplete objects; characterization principle; identity conditions; existence}

\begin{document}
\maketitle

\begin{abstract}
Meinong's jungle is the domain of objects that have properties without existing: the golden mountain, the round square, the possible fat man in the doorway. MET-24-a argues that existence is a first-order property had by everything and that there are no Meinongian objects. We add an argument that turns on a feature of the jungle's inhabitants that MET-24-a does not exploit: their incompleteness. A Meinongian object characterized by an incomplete condition neither has nor lacks the properties the condition leaves open. We show that, with plain predication and classical logic, nothing is incomplete; that completing the characterized objects destroys the uniqueness on which Meinongian descriptions rely; and that the identity conditions Meinongians supply hold only of objects that, on their own account, exist as abstract objects. Every inhabitant of the jungle either leaves the domain or exists.
\end{abstract}

\section{The inhabitants}

\citet{Routley1980} named the region of nonexistent objects Meinong's jungle. Meinong held that the object of every thought has a being-so independent of its being, so that the golden mountain is golden and a mountain though there is no such mountain \citep{Meinong1904}. Its inhabitants are of several kinds: incomplete objects such as the golden mountain, which is silent about its height; impossible objects such as the round square; merely possible objects such as Quine's possible fat man in the doorway \citep{Quine1948}; and the objects of thought and of fiction.

MET-24-a deals with most of them. Its Theorem~4.6 shows that a Meinongian who posits a golden mountain in the domain of everything uses \enquote{golden} and \enquote{mountain} in a sense other than the plain one, on pain of contradicting the plain fact that there are no golden mountains, and that the object so posited exists; such objects, MET-24-a argues, are existing abstract objects. Its Theorems~5.3 and~5.4 show that Meinongian theories are intertranslatable with classical theories in which the Meinongian's existence predicate is an ordinary predicate. Its Section~6 treats negative existentials, fiction, intentionality, merely possible and impossible objects. We add an argument from the incompleteness of characterized objects, which shows that the jungle cannot be stocked as Meinong stocked it even before the question of existence arises.

\section{Incompleteness}

\begin{definition}[Incompleteness]\label{d:incomplete}
An object $o$ is \emph{incomplete} with respect to a property $F$ if neither $Fo$ nor not-$Fo$.
\end{definition}

The characterization principle, in its nuclear form \citep{Parsons1980}, assigns to each set of nuclear properties an object that has exactly those nuclear properties. The object assigned to $\{$golden, mountain$\}$ has no height property: it is not over a thousand metres tall, and it is not the case that it is over a thousand metres tall. Meinongians accept this, and keep classical logic by distinguishing a negation internal to characterization, on which the golden mountain fails to have the height property among its characterizing properties, from the classical negation of the sentence.

\begin{proposition}[Plain completeness]\label{p:complete}
Under plain predication with classical negation, no object is incomplete with respect to any property.
\end{proposition}

\begin{proof}
For any object $o$ and property $F$, either $Fo$ or not-$Fo$, by excluded middle. Under plain predication, not-$Fo$ says that $o$ lacks $F$, which is what Definition~\ref{d:incomplete} denies.
\end{proof}

Proposition~\ref{p:complete} sets a dilemma. Either the golden mountain is incomplete, and then its predication is not plain, so that, by Theorem~4.6 of MET-24-a, it is not a mountain in the sense in which there are no golden mountains, and it exists, as an abstract object; or it is complete, and then it has a determinate height, mass and location.

\begin{proposition}[Completion destroys uniqueness]\label{p:unique}
If the characterization principle supplies complete objects, then for every condition that leaves some property open and is satisfied by any of them, it is satisfied by more than one, and the description \enquote{the object satisfying it} is improper.
\end{proposition}

\begin{proof}
Let $\varphi$ leave $F$ open, so that $\varphi \wedge F$ and $\varphi \wedge \neg F$ are both consistent with whatever restrictions the theory imposes. A principle that supplies complete objects for these conditions supplies an object satisfying $\varphi$ and $F$ and one satisfying $\varphi$ and not $F$; both satisfy $\varphi$, and they differ in $F$, so they are two.
\end{proof}

The second horn is therefore no refuge. The Meinongian relies on descriptions such as \enquote{the golden mountain} and \enquote{the round square} picking out unique objects, and on these objects having exactly the properties their descriptions state. Complete objects satisfy the second requirement only by having many further properties, and the first not at all: there are as many golden mountains as completions, and \enquote{the golden mountain} denotes none of them. The thought of the golden mountain, which was to be about a single object, is about no object in particular, which is the classical view of MET-24-a with extra objects that do no work.

\section{Identity}

\citet{Quine1948} asked how many possible fat men there are in the doorway, and whether the possible fat man is the same as the possible bald man. The question tests whether the jungle's inhabitants have identity conditions, which they must have if they are objects at all (Principle~P1 of MET-13-a). The nuclear theory has an answer: objects are identical if and only if they have the same nuclear properties. But the answer is available only for incomplete objects, since it individuates them by their characterizing sets, and Proposition~\ref{p:complete} shows that incomplete objects are objects of non-plain predication.

\begin{corollary}[Identity requires abstractness]\label{c:identity}
The identity conditions that nuclear Meinongianism supplies hold of objects whose predication is not plain, and these objects, if they are in the domain, are existing abstract objects.
\end{corollary}

\begin{proof}
By Proposition~\ref{p:complete}, an object individuated by a characterizing set of properties and incomplete outside it is not a subject of plain predication. By Theorem~4.6 of MET-24-a, an object that is in the domain and is golden and a mountain in the Meinongian's sense is so only in a non-plain sense, and exists. What has properties in such a sense, and is individuated by the set of them, is an abstract object, such as a set of properties or an abstract object that encodes them \citep{Zalta1983}.
\end{proof}

Quine's question therefore has a clean answer. Read with plain predication, there are no possible fat men in the doorway, since there are no fat men in the doorway; read with non-plain predication, there is one object for each set of properties, as many as there are sets, and each exists as an abstract object. The jungle's identity conditions are the identity conditions of sets.

\section{Objections}

\begin{objection}[Noneism needs no characterization principle]
\citet{Priest2005} does not hold that nonexistent objects have their characterizing properties at the actual world, but only at the worlds that realize how they are represented. Incompleteness at the actual world is then harmless.
\end{objection}

\begin{reply}
At the actual world, the noneist's golden mountain has neither the property of being golden nor that of being a mountain, and is complete with respect to plain predication: it is not golden. What it has is a relation to worlds at which something is golden and a mountain, and these, on TSM-07, are sets of propositions, which exist. Proposition~\ref{p:complete} is then satisfied, and the content of the thought is modelled by existing objects. What remains of noneism is the claim that the actual-world object stands in that relation without existing, and Theorem~5.3 of MET-24-a shows that this claim can be restated without loss with an ordinary predicate in place of \enquote{exists}.
\end{reply}

\begin{objection}[Existence as a sparse property]
The jungle seems overgrown only if existence is a thin property everything has. If existence is a substantive universal that some objects lack, the jungle is simply the class of objects without it.
\end{objection}

\begin{reply}
On MET-24-a, existence is constructed from identity by quantification and is not among the sparse universals of MET-02; it is an abundant property that adds nothing to what has it. The present argument does not depend on this. Propositions~\ref{p:complete} and~\ref{p:unique} concern predication and uniqueness and make no use of the nature of existence, so even a substantive existence property would leave the jungle's inhabitants either incomplete, and so not plainly golden, or complete, and so not unique.
\end{reply}

\section{Conclusion}

Meinong stocked his jungle by characterization: for each condition, an object with exactly the properties it states. With plain predication and classical logic, no object has exactly the properties of an incomplete condition. Incomplete objects are therefore objects of non-plain predication, and these exist as abstract objects; complete ones are many where the Meinongian needs one. The identity conditions Meinongians supply are those of sets of properties. Nothing in the jungle is a nonexistent object: what was there either is not in the domain or exists.

\begin{thebibliography}{99}\small
\bibitem[Meinong(1904)]{Meinong1904} Meinong, A. (1904). \"Uber Gegenstandstheorie. In A.~Meinong (ed.), \emph{Untersuchungen zur Gegenstandstheorie und Psychologie}. Barth.
\bibitem[Parsons(1980)]{Parsons1980} Parsons, T. (1980). \emph{Nonexistent Objects}. Yale University Press.
\bibitem[Priest(2005)]{Priest2005} Priest, G. (2005). \emph{Towards Non-Being: The Logic and Metaphysics of Intentionality}. Oxford University Press.
\bibitem[Quine(1948)]{Quine1948} Quine, W.~V.~O. (1948). On what there is. \emph{The Review of Metaphysics}, 2(5), 21--38.
\bibitem[Routley(1980)]{Routley1980} Routley, R. (1980). \emph{Exploring Meinong's Jungle and Beyond}. Research School of Social Sciences, Australian National University.
\bibitem[Zalta(1983)]{Zalta1983} Zalta, E.~N. (1983). \emph{Abstract Objects: An Introduction to Axiomatic Metaphysics}. Reidel.
\end{thebibliography}

\end{document}
